\chapter{Compact Operators and Fredholm Theory}
\label{chap:compact}

Serre's \emph{Endomorphismes compl\`etement continus des espaces de Banach $p$-adiques}
\cite{Serre} attaches to a completely continuous operator $u$ on a $p$-adic Banach space a
characteristic power series $\det(1 - Tu)$, entire in $T$, whose reciprocal roots are exactly the
eigenvalues of $u$, and which controls a Riesz decomposition of the space.  That is the machine
we need: a Hecke operator will be completely continuous, its determinant will be computable from
a matrix, and the Newton polygon of the determinant will read off the eigenvalue valuations.

Serre works over a complete nonarchimedean field.  Buzzard \cite{Buzzard} works over a Banach
algebra over such a field, and Johansson--Newton \cite{JN} over a Banach--Tate ring with no
ground field at all.  We formalise the last of these, since the applications need the ground
ring to be an affinoid rather than a field, and since nothing in the arguments actually uses a
field.  The specialisations are then corollaries.

\section{The setting}

\begin{definition}
  \label{def:istate}
  \lean{TateFredholm.IsTate}
  \leanok
  A \emph{Banach--Tate ring} is a complete normed ring $R$ possessing a
  \emph{pseudo-uniformizer}: a unit $\varpi$ with $\lVert \varpi \rVert < 1$ whose norm is
  multiplicative.  We write $\mathrm{IsTate}\, R$.
\end{definition}

Note that the norm on $R$ is only submultiplicative; the pseudo-uniformizer is the one element
whose norm behaves multiplicatively, and that turns out to be all the analysis needs.

\begin{definition}
  \label{def:cspace}
  \lean{TateFredholm.cSpace}
  \leanok
  For an index type $I$ the \emph{model space} $c(I, R)$ is the space of families $(x_i)_{i \in
  I}$ in $R$ tending to $0$ along the cofinite filter, with the supremum norm.
\end{definition}

Every module we care about is either $c(I,R)$ or a direct summand of one:

\begin{definition}
  \label{def:onable}
  \uses{def:cspace}
  \lean{TateFredholm.IsONable, TateFredholm.HasPr}
  \leanok
  A Banach $R$-module $M$ is \emph{orthonormalisable} if it is isometrically isomorphic to some
  $c(I,R)$, and has \emph{property (Pr)} if it is a direct summand of an orthonormalisable
  module.
\end{definition}

Property (Pr) is Buzzard's hypothesis, and it is what a space of modular forms at a non-neat
level satisfies when orthonormalisability itself fails.

\begin{proposition}
  \label{prop:tate-algebra-onable}
  \uses{def:onable, def:tate-algebra}
  \lean{TateFredholm.restrictedEquivCSpace, TateFredholm.isONable_restricted}
  \leanok
  The Tate algebra $K\langle X \rangle$ is orthonormalisable, with the monomials $X^n$ as
  orthonormal basis; the isometry $K\langle X\rangle \cong c(\N, K)$ sends a series to its
  coefficient family.
\end{proposition}

This proposition is the reason the weight-$\kappa$ action of Chapter~\ref{chap:qmf}, which acts
on a Tate algebra, is an operator on a model space and so has matrix coefficients.

\section{Compactness}

\begin{definition}
  \label{def:matrixcoeff}
  \uses{def:cspace}
  \lean{TateFredholm.matrixCoeff, TateFredholm.rowNorm}
  \leanok
  For $u : c(I,R) \to c(J,R)$ continuous linear, the \emph{matrix coefficient} $u_{j i}$ is the
  $j$-th coordinate of $u(e_i)$, where $e_i$ is the $i$-th basis vector.  The $j$-th
  \emph{row norm} is $\sup_i \lVert u_{ji}\rVert$.
\end{definition}

\begin{definition}
  \label{def:iscompactoid}
  \uses{def:matrixcoeff}
  \lean{TateFredholm.IsCompletelyContinuous, TateFredholm.IsCompactoid}
  \leanok
  $u$ is \emph{completely continuous} if it is an operator-norm limit of finite-rank operators,
  and \emph{compactoid} if its row norms tend to $0$ along the cofinite filter.
\end{definition}

\begin{proposition}
  \label{prop:compactoid-criterion}
  \uses{def:iscompactoid}
  \lean{TateFredholm.isCompactoid_of_row_decay'}
  \leanok
  On $c(I,R)$ the two notions agree, and a row-decay estimate $\lVert u_{ji} \rVert \leq C
  \sigma^{j}$ with $\sigma < 1$ suffices for both.
\end{proposition}

The row-decay criterion is the only form of compactness we ever verify in practice: for a Hecke
operator the estimate falls out of the weight action, see Theorem~\ref{thm:qmf-compact}.

\section{The Fredholm determinant}

\begin{definition}
  \label{def:minor}
  \uses{def:matrixcoeff}
  \lean{TateFredholm.minor, TateFredholm.charCoeff, TateFredholm.charPowerSeries}
  \leanok
  For a finite $S \subseteq I$ let $\mathrm{minor}(u, S)$ be the determinant of the $S \times S$
  submatrix of $u$.  For compactoid $u$ the family of $n$-element minors is summable, and
  \[
    c_n(u) = (-1)^n \sum_{\lvert S \rvert = n} \mathrm{minor}(u, S), \qquad
    \det(1 - Tu) = \sum_{n \geq 0} c_n(u) T^n .
  \]
\end{definition}

\begin{lemma}[Ultrametric Hadamard bound]
  \label{lem:hadamard}
  \uses{def:minor}
  \lean{TateFredholm.norm_det_le_pow_of_row_bound}
  \leanok
  If every entry of a square matrix satisfies $\lVert A_{ji}\rVert \leq \sigma^{w_j}$ for a row
  weight $w$, then $\lVert \det A \rVert \leq \sigma^{\sum_j w_j}$.
\end{lemma}

\begin{proof}
  \leanok
  Each Leibniz monomial picks one entry from every row, so has norm at most $\prod_j
  \sigma^{w_j}$; the ultrametric inequality bounds the sum of the monomials by their maximum.
\end{proof}

Summability of the minors comes from the same bound: a minor of norm $\geq \varepsilon$ forces
its index set inside the finite set of ``large'' rows, and there are finitely many $n$-subsets of
a finite set.

\begin{theorem}
  \label{thm:charpowerseries-entire}
  \uses{def:minor}
  \lean{TateFredholm.charPowerSeries_isEntire, TateFredholm.charCoeff_zero}
  \leanok
  For compactoid $u$ the series $\det(1 - Tu)$ is entire, with constant term $1$.
\end{theorem}

\begin{theorem}[Invariances]
  \label{thm:charpowerseries-inv}
  \uses{def:minor}
  \lean{TateFredholm.charPowerSeries_comm, TateFredholm.charPowerSeries_conj}
  \leanok
  $\det(1 - Tuv) = \det(1 - Tvu)$, and the determinant is invariant under conjugation by an
  isomorphism.  Consequently it may be computed in any model of the space.
\end{theorem}

\begin{proposition}[Block operators]
  \label{prop:blockop}
  \uses{def:minor}
  \lean{TateFredholm.blockOp, TateFredholm.isCompactoid_blockOp}
  \leanok
  An operator on $c(\sigma \times I, R)$ with $\sigma$ finite is assembled from its $\sigma
  \times \sigma$ blocks; it is compactoid iff each block is, and its determinant is computed
  blockwise.
\end{proposition}

This is the shape a Hecke operator takes: the class set indexes $\sigma$, and each block is the
weight action of a single certificate element.

\section{Serre's theorems}

\begin{theorem}[{\cite[Prop.\ 11]{Serre}}]
  \label{thm:serre11}
  \uses{thm:charpowerseries-entire}
  \lean{TateFredholm.isUnit_one_sub_smul_iff_isUnit_evalT}
  \leanok
  For compactoid $u$ and $a \in R$, the operator $1 - a u$ is invertible if and only if
  $\det(1-Tu)$ evaluated at $a$ is a unit.
\end{theorem}

\begin{theorem}[{\cite[Prop.\ 12]{Serre}}]
  \label{thm:serre12}
  \uses{thm:serre11}
  \lean{TateFredholm.exists_eigenvector_of_evalT_charPowerSeries_eq_zero}
  \leanok
  Over a field, if $a \neq 0$ is a zero of $\det(1 - Tu)$ then $a^{-1}$ is an eigenvalue of $u$;
  and conversely.
\end{theorem}

Theorem~\ref{thm:serre12} is the statement that makes the Newton polygon of the determinant say
something about the operator, and hence the hinge of the whole thesis.

\begin{theorem}[Riesz decomposition]
  \label{thm:riesz}
  \uses{thm:serre12}
  \lean{TateFredholm.exists_riesz_decomposition}
  \leanok
  Let $u$ be compactoid on $c(I,K)$ over a discretely valued complete field, and let $a$ be a
  zero of $\det(1 - Tu)$ of multiplicity $h$.  Then $c(I,K) = N \oplus F$ with both summands
  $u$-stable, $(1 - au)^h$ vanishing on $N$, $(1-au)$ invertible on $F$, and
  \[
    \dim_K N = h .
  \]
\end{theorem}

The subspace $N$ is the \emph{finite-slope subspace} attached to $a$; it is finite dimensional,
which is what makes eigenvalue statements about it meaningful, and everything commuting with $u$
preserves it, which is what makes a Hecke algebra act on it (Proposition~\ref{prop:qmf-eigensystem}).

%% The projector p = e^h and Serre's Lemme 2 (block-triangular determinants) are the internal
%% steps: TateFredholm.exists_rieszProjection, charPowerSeries_blockTriangular.
